\documentclass[11pt,a4paper]{article}

% --- Core LaTeX Packages ---
\usepackage[utf8]{inputenc}
\usepackage[margin=1in]{geometry}
\usepackage{amsmath,amssymb,amsthm}
\usepackage{xcolor}
\usepackage{listings}
\usepackage{hyperref}
\usepackage{tabularx}
\usepackage{graphicx}

% --- Optional Diagram Packages (include when needed) ---
\usepackage{forest}
% \usepackage{tikz}

% --- Unified Color Palette ---
\definecolor{primaryblue}{RGB}{24, 75, 140}
\definecolor{codebg}{RGB}{248, 249, 250}
\definecolor{codeframe}{RGB}{210, 215, 222}
\definecolor{codegreen}{RGB}{40, 130, 60}
\definecolor{codeblue}{RGB}{20, 90, 180}
\definecolor{codegray}{RGB}{120, 120, 120}

% --- Hyperref Setup ---
\hypersetup{
    colorlinks=true,
    linkcolor=primaryblue,
    citecolor=primaryblue,
    urlcolor=primaryblue,
    pdftitle={Recursion and Backtracking},
    pdfauthor={Antigravity Engineering}
}

% --- Theorem & Definition Environments ---
\theoremstyle{definition}
\newtheorem{definition}{Definition}[section]
\newtheorem{theorem}{Theorem}[section]
\newtheorem{lemma}[theorem]{Lemma}
\newtheorem{example}{Example}[section]

% --- Callout Box Environment ---
\newenvironment{calloutbox}[1]{
    \par\vspace{0.3cm}
    \noindent\begin{tabular}{|p{0.96\textwidth}|}
    \hline
    \vspace{0.1cm}
    \textbf{#1}\par\vspace{0.1cm}
}{
    \vspace{0.1cm}
    \\ \hline
    \end{tabular}
    \vspace{0.3cm}\par
}

% --- Modern C++ Code Listing Setup ---
\lstset{
    language=C++,
    basicstyle=\footnotesize\ttfamily,
    keywordstyle=\color{codeblue}\bfseries,
    commentstyle=\color{codegreen}\itshape,
    stringstyle=\color{orange!80!black},
    numberstyle=\tiny\color{codegray},
    numbers=left,
    numbersep=8pt,
    backgroundcolor=\color{codebg},
    frame=single,
    rulecolor=\color{codeframe},
    breaklines=true,
    breakatwhitespace=true,
    tabsize=4,
    showstringspaces=false,
    captionpos=b
}

% --- Title Block ---
\title{\vspace{-1.5cm}\textbf{\Huge Recursion and Backtracking}\\[0.3cm]
\Large Call Stacks, Recursive Trees, Pruning Strategies, and C++ Implementations}
\author{Technical Monograph}
\date{\today}

\begin{document}

\maketitle

\begin{abstract}
\noindent
This document explains recursion and backtracking from the ground up. It starts with the absolute basics of function calls and the call stack, explaining how a function can call itself to solve a smaller piece of a problem. Then, it introduces backtracking as a way to systematically explore all possible choices to find valid solutions, like navigating a maze. Each concept has a plain-English explanation, step-by-step traces of how the algorithms explore choices, and practical C++ implementations. The final sections help you look at new problems and decide exactly when and how to apply these techniques.
\end{abstract}

\tableofcontents
\newpage

\section{Foundations: The Call Stack and Recursion}

\subsection{The Call Stack}

Before you can understand recursion, you have to understand what happens when a function calls another function. Imagine your computer has a stack of blank sticky notes. Every time a function is called, the computer takes a new sticky note, writes down the function's variables and what line of code it is currently on, and pushes this note onto the top of the stack. This note is called a \textbf{stack frame}.

While that top function is running, all the other functions underneath it are paused. If the top function calls yet another function, a new sticky note is added on top. When a function finishes running, its sticky note is thrown away (popped off the stack), and the computer resumes whatever function is now on top, right where it left off.

\begin{calloutbox}{Stack Overflow Warning}
Every recursive call uses stack memory for its frame. The default stack size is typically 1--8 MB. For recursion depth $n$, if each frame uses $f$ bytes, you need $n \cdot f$ bytes of stack. For $n = 10,000$ with frames of 100 bytes, that is 1 MB --- borderline. For very deep recursion ($n > 10,000$), prefer an iterative solution with an explicit stack.
\end{calloutbox}

\subsection{Recursion}

\textbf{Recursion} is simply what happens when a function calls itself. It is a way of solving a large problem by breaking it down into smaller and smaller pieces of the exact same problem, until the pieces are so small that they are easy to solve directly. 

Every recursive process needs a \textbf{base case}. This is the sticky note that says ``stop.'' Without a base case, the function would keep calling itself forever until it runs out of sticky notes (a stack overflow).

\textbf{Simplest Example: Factorial of 5}
The factorial of 5 (written as $5!$) is $5 \times 4 \times 3 \times 2 \times 1$. 
Observe that $5!$ is exactly the same as $5 \times 4!$. And $4!$ is exactly $4 \times 3!$. 

Here is how the computer computes $5!$ step by step using recursion:
\begin{itemize}
    \item \textbf{Call factorial(5):} Is 5 the base case? No. Return $5 \times \text{factorial}(4)$. (Pushes sticky note for 5).
    \item \textbf{Call factorial(4):} Is 4 the base case? No. Return $4 \times \text{factorial}(3)$. (Pushes sticky note for 4).
    \item \textbf{Call factorial(3):} Is 3 the base case? No. Return $3 \times \text{factorial}(2)$. (Pushes sticky note for 3).
    \item \textbf{Call factorial(2):} Is 2 the base case? No. Return $2 \times \text{factorial}(1)$. (Pushes sticky note for 2).
    \item \textbf{Call factorial(1):} 1 is the base case! Return 1. (No new sticky note. Start unwinding).
\end{itemize}
Now the stack unwinds, multiplying the returned values back up the chain: $1 \times 2 = 2$, $2 \times 3 = 6$, $6 \times 4 = 24$, $24 \times 5 = 120$.

\subsection{Recursion Tree}

A \textbf{recursion tree} is a visual representation of all the recursive calls made by an algorithm. When a function calls itself multiple times, the tree branches out. 

\textbf{Worked Example: Fibonacci(4)}
The Fibonacci sequence defines the $n$-th number as the sum of the previous two numbers: $\text{fib}(n) = \text{fib}(n-1) + \text{fib}(n-2)$, with base cases $\text{fib}(0) = 0$ and $\text{fib}(1) = 1$. Let us draw the recursion tree for $\text{fib}(4)$:

Observe how many times sub-problems are re-computed. The call \texttt{fib(2)} is computed entirely from scratch twice. \texttt{fib(1)} is computed three times. 

\subsection{The Three Laws of Recursion}
For a recursive algorithm to be correct, it must obey three rules:
\begin{enumerate}
    \item A recursive algorithm must have a base case.
    \item A recursive algorithm must change its state and move toward the base case.
    \item A recursive algorithm must call itself, recursively.
\end{enumerate}

\section{Backtracking: Choose, Explore, Undo}

\subsection{Backtracking}

\textbf{Backtracking} is a systematic way to search for solutions to a problem by exploring all possible options, but abandoning a path as soon as you realize it cannot lead to a valid solution. 

Think of it as navigating a maze. You reach an intersection with three paths. You pick the first path and walk down it. If it leads to a dead end, you walk back to the intersection (\textbf{undo the choice}) and try the next path. You are doing a depth-first search through a \textbf{decision tree}.

\begin{calloutbox}{Backtracking vs. Brute Force}
Brute force generates every possible configuration and then filters. Backtracking checks constraints incrementally and abandons a branch the moment it becomes invalid. For generate parentheses, brute force generates all $2^{2n}$ binary strings of length $2n$ and then checks validity --- for $n=8$ that is 65,536 strings. Backtracking explores only the $C_8 = 1430$ valid ones plus a few pruned branches.
\end{calloutbox}

\subsection{The Backtracking Template}

The standard pattern for all backtracking problems looks like this:

\begin{calloutbox}{The One Pattern Behind All Backtracking}
The make-a-choice $\rightarrow$ recurse $\rightarrow$ undo-the-choice structure is always the same. The only thing that changes between different backtracking problems is what ``choice'' means and what the constraints are.
\end{calloutbox}

\subsection{Pruning}

\textbf{Pruning} is the act of stopping the exploration of a branch early because you know it is hopeless. There are two types of rules for pruning:
\begin{itemize}
    \item \textbf{Constraint}: A hard rule. Violating it makes the path invalid. Pruning on hard constraints is always safe and always beneficial.
    \item \textbf{Optimality Condition}: A soft rule. The path might be valid, but you already have a better solution, so it is suboptimal. 
\end{itemize}

\subsection{Generate Parentheses: Constraint-Based Pruning}

Consider generating all valid combinations of $n$ pairs of parentheses. Let us use $n=2$. We have 4 slots to fill. At each slot, we can place a \texttt{(} or a \texttt{)}. 

Our two hard constraints for pruning:
\begin{enumerate}
    \item The total number of open parentheses cannot exceed $n$.
    \item The total number of close parentheses cannot exceed the number of open parentheses placed so far (otherwise, we have an unmatched closing bracket).
\end{enumerate}

Step-by-step trace for $n=2$:
\begin{itemize}
    \item Start: \texttt{""} (open: 0, close: 0).
    \item \textbf{Choice \texttt{(}}: \texttt{(} (open: 1, close: 0). Valid.
        \begin{itemize}
            \item \textbf{Choice \texttt{(}}: \texttt{((} (open: 2, close: 0). Valid.
                \begin{itemize}
                    \item \textbf{Choice \texttt{(}}: Cannot place open, open == 2. \textbf{PRUNED}.
                    \item \textbf{Choice \texttt{)}}: \texttt{(()} (open: 2, close: 1). Valid.
                        \begin{itemize}
                            \item \textbf{Choice \texttt{)}}: \texttt{(())} (open: 2, close: 2). Base case! Record solution.
                        \end{itemize}
                \end{itemize}
            \item \textbf{Choice \texttt{)}}: \texttt{()} (open: 1, close: 1). Valid.
                \begin{itemize}
                    \item \textbf{Choice \texttt{(}}: \texttt{()(} (open: 2, close: 1). Valid.
                        \begin{itemize}
                            \item \textbf{Choice \texttt{)}}: \texttt{()()} (open: 2, close: 2). Base case! Record solution.
                        \end{itemize}
                    \item \textbf{Choice \texttt{)}}: Cannot place close, close == open. \textbf{PRUNED}.
                \end{itemize}
        \end{itemize}
    \item \textbf{Choice \texttt{)}}: Cannot place close, close == open. \textbf{PRUNED}.
\end{itemize}

\subsection{Subsets, Permutations, and Combinations}

\textbf{Subsets (Power Set):} Generating all subsets of $\{1, 2, 3\}$. The decision at each step is simply: do we include the current element, or exclude it?
For 3 elements, we make 3 yes/no decisions, leading to $2^3 = 8$ subsets.

\textbf{Permutations:} Generating all orderings of $[1, 2, 3]$. We want arrays of length 3 where each element appears exactly once. Our choice is picking which available number to place next. We use a \texttt{used} array to keep track of which numbers we have already chosen. Where does this come from? Intuition: If we do not track what is used, we might pick the exact same index twice, which would just give us length-3 combinations with replacement. The \texttt{used} array effectively shrinks the available pool without needing to resize the underlying container.

\textbf{Combinations:} Choosing $k$ elements from $n$. To avoid generating duplicates (like picking [1,2] and then [2,1]), we enforce an order. We pass a \texttt{start} index into our recursive function so that we only pick elements that appear \textit{after} our last pick. Where does this come from? Intuition: Combinations do not care about order, so we can enforce that every valid combination is strictly sorted. The \texttt{start} index guarantees we never pick an element smaller than what we already picked.

\section{Key Algorithmic Patterns}

\textbf{Why the naive approach fails:}
Imagine you want to find all valid permutations of 10 elements. The naive approach is to generate all sequences and filter the invalid ones. There are $10! = 3,628,800$ combinations. With 12 elements, it jumps to 479 million. If you have constraints, many branches can be pruned before they are ever completed. If you know the first element being \texttt{5} invalidates the puzzle, backtracking skips the $9! = 362,880$ permutations that start with \texttt{5}. Brute force generates all of them anyway.

\textbf{The Efficient Technique:}
The core pattern is treating the \textbf{search space as a tree, performing a depth-first traversal (DFS), and pruning early}. You maintain a single piece of state (like a string or an array). As you go deeper down the tree, you modify this state. When you return from the recursive call (backtrack), you revert the state to exactly how it was before. This allows you to explore massive search spaces with extremely minimal memory (just the call stack and one state variable).

\textbf{Step-by-Step Trace: Generate Parentheses for n=2}
\begin{enumerate}
    \item State is empty string \texttt{""}.
    \item Try adding \texttt{(}. State is \texttt{"("}. Recurse.
    \item Try adding \texttt{(}. State is \texttt{"(("}. Recurse.
    \item Try adding \texttt{)}. State is \texttt{"(()"}. Recurse.
    \item Try adding \texttt{)}. State is \texttt{"(())"}. Length is 4. Add to results. Return.
    \item Backtrack: undo \texttt{)} to get \texttt{"(()"}. 
    \item Backtrack: undo \texttt{)} to get \texttt{"(("}. 
    \item Backtrack: undo \texttt{(} to get \texttt{"("}.
    \item Try adding \texttt{)}. State is \texttt{"()"}. Recurse.
    \item Try adding \texttt{(}. State is \texttt{"()("}. Recurse.
    \item Try adding \texttt{)}. State is \texttt{"()()"}. Length is 4. Add to results. Return.
\end{enumerate}

\subsection{The Make-a-Choice $\rightarrow$ Explore $\rightarrow$ Undo Template}

Backtracking is elegant because you do not have to copy massive states. You share one global state (like a \texttt{current} list). You append to it, recursively explore, and then pop the last element off to restore it for the next choice.

Where does this come from? The intuition is that each recursive function invocation conceptually represents a node in the decision tree. To visit its sibling node, you must reverse whatever side-effects the current node caused on the shared state, returning it to the exact condition it was in when the parent node handed it down.

\subsection{Handling Duplicate Elements}
When the input array contains duplicate elements (for example, \texttt{[1, 2, 2]}), standard backtracking will generate duplicate subsets or permutations.
The efficient pattern to handle this is:
\begin{enumerate}
    \item \textbf{Sort the input} so that duplicate elements are adjacent.
    \item \textbf{Skip equal siblings}: Before making a choice in your loop, check if the current element is equal to the previous element. If it is, and the previous element was skipped (or not used at this level), skip the current element too.
\end{enumerate}

For subsets of \texttt{[1, 2, 2]}:
\begin{itemize}
    \item Pick \texttt{1} $\rightarrow$ Pick first \texttt{2} $\rightarrow$ \texttt{[1, 2]}.
    \item Backtrack to \texttt{1}. Next choice is the second \texttt{2}.
    \item Since the second \texttt{2} is identical to the first \texttt{2} (which we just finished exploring), we skip it to avoid generating another \texttt{[1, 2]}.
\end{itemize}

\section{Worked Examples}

\subsection{Example 1: Factorial}
Compute $4!$.
\begin{enumerate}
    \item call \texttt{fac(4)}. Base case? No. Returns \texttt{4 * fac(3)}.
    \item call \texttt{fac(3)}. Base case? No. Returns \texttt{3 * fac(2)}.
    \item call \texttt{fac(2)}. Base case? No. Returns \texttt{2 * fac(1)}.
    \item call \texttt{fac(1)}. Base case? Yes. Returns \texttt{1}.
    \item \texttt{fac(2)} receives 1, computes $2 \times 1 = 2$, returns 2.
    \item \texttt{fac(3)} receives 2, computes $3 \times 2 = 6$, returns 6.
    \item \texttt{fac(4)} receives 6, computes $4 \times 6 = 24$, returns 24.
\end{enumerate}

\subsection{Example 2: Generate Parentheses (n=3)}
The total number of valid strings for $n$ pairs is given by the $n$-th \textbf{Catalan number}, $C_n$. For $n=3$, $C_3 = 5$. The valid configurations are:
1. \texttt{((()))} \\
2. \texttt{(()())} \\
3. \texttt{(())()} \\
4. \texttt{()(())} \\
5. \texttt{()()()}

\vspace{0.5cm}
\noindent\textbf{Decision Tree for n=2}:

\subsection{Example 3: Subsets of \{1, 2\}}
We start with an empty set. At each element, we can include or exclude. 
\begin{enumerate}
    \item Exclude 1 $\rightarrow$ Exclude 2 $\rightarrow$ \texttt{[]}
    \item Exclude 1 $\rightarrow$ Include 2 $\rightarrow$ \texttt{[2]}
    \item Include 1 $\rightarrow$ Exclude 2 $\rightarrow$ \texttt{[1]}
    \item Include 1 $\rightarrow$ Include 2 $\rightarrow$ \texttt{[1, 2]}
\end{enumerate}

\vspace{0.5cm}
\noindent\textbf{Decision Tree for Subsets of \{1, 2\}}:

\subsection{Example 4: Permutations of [1, 2, 3]}
We want all orderings.
\begin{enumerate}
    \item Start empty. Choose 1. Remaining: 2, 3.
    \item Choose 2. Remaining: 3. Choose 3 $\rightarrow$ \texttt{[1, 2, 3]}.
    \item Backtrack to when remaining was 2, 3. Choose 3. Remaining: 2. Choose 2 $\rightarrow$ \texttt{[1, 3, 2]}.
    \item Backtrack to start. Choose 2. Remaining: 1, 3.
    \item Choose 1 $\rightarrow$ \texttt{[2, 1, 3]}. Choose 3 $\rightarrow$ \texttt{[2, 3, 1]}.
    \item Backtrack to start. Choose 3. Remaining: 1, 2.
    \item Choose 1 $\rightarrow$ \texttt{[3, 1, 2]}. Choose 2 $\rightarrow$ \texttt{[3, 2, 1]}.
\end{enumerate}

\subsection{Example 5: N-Queens (N=4)}
Place 4 queens on a $4 \times 4$ board so none attack each other (no two in the same row, column, or diagonal). Let us trace it and see pruning in action. We place row by row.
\begin{enumerate}
    \item Row 0: Place queen at column 0.
    \item Row 1: Columns 0 and 1 are attacked. Place at column 2.
    \item Row 2: Columns 0, 1, 2, 3 are all attacked! Dead end. \textbf{PRUNED}. Backtrack to Row 1.
    \item Row 1: Move queen to column 3.
    \item Row 2: Place at column 1.
    \item Row 3: All columns attacked. Dead end. \textbf{PRUNED}. Backtrack.
    \item Backtrack all the way to Row 0. Move queen to column 1.
    \item Row 1: Place at column 3.
    \item Row 2: Place at column 0.
    \item Row 3: Place at column 2. Valid! (Solution 1)
\end{enumerate}

By checking constraints before diving into deeper rows, we avoid generating thousands of invalid boards. Total solutions for N=4 is 2. (For 8-Queens it is 92).

\section{C++ Implementations}

\subsection{Simple Recursion}
This is a straightforward recursive function with no choices and no backtracking.

\begin{lstlisting}[caption={Factorial using recursion}]
#include <cstdint>

[[nodiscard]] constexpr int64_t factorial(int n) noexcept {
    // 1. Base case: 0! and 1! are 1
    if (n <= 1) {
        return 1;
    }
    // 2. Recursive step: n * (n-1)!
    return n * factorial(n - 1);
}
\end{lstlisting}

\subsection{Constraint-Based Backtracking}
At each position in the string, our choice is simple: append an open parenthesis \texttt{(}, or append a close parenthesis \texttt{)}. We only make the choice if it obeys our hard constraints (open $<$ n, close $<$ open). A note on depth: the stack depth here is $2n$, which easily fits in memory.

\begin{lstlisting}[caption={Generate Parentheses using Backtracking}]
#include <vector>
#include <string>

void backtrackParentheses(int n, int open, int close, std::string& current, std::vector<std::string>& result) {
    // 1. Base case: If we have placed 2*n characters, the string is complete.
    if (open == n && close == n) {
        result.push_back(current);
        return;
    }
    
    // 2. Make choice: Add an open parenthesis if we haven't reached the limit.
    if (open < n) {
        current.push_back('('); // Make choice
        backtrackParentheses(n, open + 1, close, current, result); // Recurse
        current.pop_back(); // Undo choice
    }
    
    // 3. Make choice: Add a close parenthesis if there are unmatched open ones.
    if (close < open) {
        current.push_back(')'); // Make choice
        backtrackParentheses(n, open, close + 1, current, result); // Recurse
        current.pop_back(); // Undo choice
    }
}

[[nodiscard]] std::vector<std::string> generateParenthesis(int n) {
    std::vector<std::string> result;
    std::string current;
    backtrackParentheses(n, 0, 0, current, result);
    return result;
}
\end{lstlisting}

\subsection{Subsets, Permutations, and Combinations}
Here are the core backtracking implementations for generating subsets, permutations, and combinations in modern C++23. Notice the use of \texttt{std::ranges::sort} to prepare data when needed.

\begin{lstlisting}[caption={Subsets (Power Set)}]
#include <vector>

void backtrackSubsets(const std::vector<int>& nums, int index, std::vector<int>& current, std::vector<std::vector<int>>& result) {
    // 1. Every path in subsets is a valid combination, so record it immediately
    result.push_back(current);
    
    // 2. Iterate through all remaining elements to include
    for (int i = index; i < (int)nums.size(); ++i) {
        current.push_back(nums[i]); // Include nums[i]
        backtrackSubsets(nums, i + 1, current, result); // Recurse moving forward
        current.pop_back(); // Exclude nums[i] (undo)
    }
}

[[nodiscard]] std::vector<std::vector<int>> subsets(const std::vector<int>& nums) {
    std::vector<std::vector<int>> result;
    std::vector<int> current;
    backtrackSubsets(nums, 0, current, result);
    return result;
}
\end{lstlisting}

\begin{lstlisting}[caption={Permutations}]
#include <vector>

void backtrackPermutations(const std::vector<int>& nums, std::vector<bool>& used, std::vector<int>& current, std::vector<std::vector<int>>& result) {
    // 1. Base case: The permutation is complete when it has all elements
    if (current.size() == nums.size()) {
        result.push_back(current);
        return;
    }
    
    // 2. Try placing each available number in the current slot
    for (int i = 0; i < (int)nums.size(); ++i) {
        if (used[i]) continue; // Skip numbers already placed
        
        used[i] = true;
        current.push_back(nums[i]);
        
        backtrackPermutations(nums, used, current, result); // Recurse
        
        current.pop_back(); // Undo choice
        used[i] = false; // Mark number as available again
    }
}

[[nodiscard]] std::vector<std::vector<int>> permutations(const std::vector<int>& nums) {
    std::vector<std::vector<int>> result;
    std::vector<int> current;
    std::vector<bool> used(nums.size(), false);
    backtrackPermutations(nums, used, current, result);
    return result;
}
\end{lstlisting}

\begin{lstlisting}[caption={Combinations (n choose k)}]
#include <vector>

void backtrackCombinations(int n, int k, int start, std::vector<int>& current, std::vector<std::vector<int>>& result) {
    // 1. Base case: We have chosen exactly k elements
    if ((int)current.size() == k) {
        result.push_back(current);
        return;
    }
    
    // 2. Pick elements from 'start' to 'n' to avoid order duplicates
    for (int i = start; i <= n; ++i) {
        current.push_back(i);
        backtrackCombinations(n, k, i + 1, current, result);
        current.pop_back();
    }
}

[[nodiscard]] std::vector<std::vector<int>> combinations(int n, int k) {
    std::vector<std::vector<int>> result;
    std::vector<int> current;
    backtrackCombinations(n, k, 1, current, result);
    return result;
}
\end{lstlisting}

\subsection{N-Queens Constraint Satisfaction}
The N-Queens problem is a classic example of constraint satisfaction. By tracking used columns and diagonals in O(1) space arrays, we vastly speed up validation over a loop. Why do diagonal sets work? Intuition: on a grid, for any cell $(r, c)$, the difference $r - c$ is constant for all cells on its primary diagonal, and the sum $r + c$ is constant for its secondary diagonal.

\begin{lstlisting}[caption={N-Queens Constraint Satisfaction with O(1) Checks}]
#include <vector>
#include <string>

void backtrackQueens(int row, int n, std::vector<std::string>& board, 
                     std::vector<bool>& cols, std::vector<bool>& diag1, std::vector<bool>& diag2, 
                     std::vector<std::vector<std::string>>& result) {
    // 1. Base case: All queens placed on all rows successfully
    if (row == n) {
        result.push_back(board);
        return;
    }
    
    // 2. Try placing a queen in each column of the current row
    for (int col = 0; col < n; ++col) {
        int d1 = row - col + n - 1; // Primary diagonal index
        int d2 = row + col;         // Secondary diagonal index
        
        if (!cols[col] && !diag1[d1] && !diag2[d2]) {
            board[row][col] = 'Q'; // Make choice
            cols[col] = diag1[d1] = diag2[d2] = true;
            
            backtrackQueens(row + 1, n, board, cols, diag1, diag2, result); // Recurse to next row
            
            board[row][col] = '.'; // Undo choice
            cols[col] = diag1[d1] = diag2[d2] = false;
        }
    }
}

[[nodiscard]] std::vector<std::vector<std::string>> solveNQueens(int n) {
    std::vector<std::vector<std::string>> result;
    std::vector<std::string> board(n, std::string(n, '.'));
    std::vector<bool> cols(n, false);
    std::vector<bool> diag1(2 * n - 1, false);
    std::vector<bool> diag2(2 * n - 1, false);
    
    backtrackQueens(0, n, board, cols, diag1, diag2, result);
    return result;
}
\end{lstlisting}

\section{Problem Pattern Recognition}

\textbf{Step 1: What kind of answer does the problem ask for?}
When faced with a combinatorics or configuration problem, your first job is to look at the required output:
\begin{itemize}
    \item ``Generate ALL valid combinations/subsets/permutations'': You must use \textbf{backtracking}. You actually have to build the strings or arrays.
    \item ``Count the number of valid configurations'': You can use backtracking to count, but \textbf{Dynamic Programming (DP)} is often vastly faster. Rule of thumb: if you need the actual configurations, use backtracking; if you only need the count, use DP or math.
    \item ``Find ONE valid configuration'': Use \textbf{backtracking with an early return}. As soon as a solution is found, stop searching other branches.
\end{itemize}

\begin{calloutbox}{When to Use DP Instead}
If the problem asks only ``how many ways?'' and not ``list all ways'', dynamic programming is almost always faster than backtracking. Backtracking enumerates every solution; DP counts them without listing them. For example, if you are asked to list all valid parenthesis strings, backtracking is correct. If you are asked only for the count, the Catalan number formula or DP would solve it much faster.
\end{calloutbox}

\textbf{Step 2: Identify the key constraint that selects the tool.}
The size of $n$ tells you everything you need to know about feasibility:
\begin{itemize}
    \item \textbf{$n \le 20$}: Standard backtracking (subsets or combinations) will usually pass without issue.
    \item \textbf{$n > 20$ for permutations}: Likely impossible to generate them all. You must either find a way to prune massively, or the problem requires a completely different approach (greedy, math).
    \item \textbf{Constraints that eliminate branches early}: If placing one piece incorrectly ruins all subsequent choices, pruning is highly effective (like N-Queens).
    \item \textbf{No constraints / all combinations valid}: Backtracking degenerates to brute force. If you just need a count, use DP or combinatorics formulas.
    \item \textbf{Input has repeated elements}: You need to sort and apply skip-equal-sibling logic to avoid duplicate outputs.
\end{itemize}

\textbf{Step 3: Is the input too large for a direct solution?}
If $n = 100$, doing $O(2^n)$ backtracking means $2^{100}$ operations, which would take longer than the age of the universe. At $n > 25$, pure exponential backtracking starts to time out in modern systems unless pruning cuts off 99\% of the tree.

\textbf{Quick Reference Table}

\noindent
\begin{tabularx}{\textwidth}{|p{4.2cm}|X|p{4cm}|}
\hline
\textbf{Keywords in the Problem} & \textbf{Plain English Meaning} & \textbf{Tool to Use} \\ \hline
``Return all valid parenthesis...'' & Build all correct bracket strings & Backtracking \\ \hline
``Find all subsets'', ``power set'' & For every element, yes or no & Backtracking (Subsets) \\ \hline
``Return all possible permutations'' & Find all distinct orderings & Backtracking (with \texttt{used} array) \\ \hline
``Return all combinations of k'' & Pick exactly $k$ out of $n$ without order & Backtracking (with \texttt{start} index) \\ \hline
``Solve the constraint puzzle / N-Queens'' & Constraint satisfaction puzzle & Backtracking (Pruning) \\ \hline
``How many distinct ways...'' & Return just the integer count & Dynamic Programming / Math \\ \hline
``Find any valid ordering...'' & Stop when you find one success & Backtracking (Early exit) \\ \hline
``Unique permutations'', ``contains duplicates'' & Find orderings without creating duplicates & Backtracking (Sort + skip identical elements) \\ \hline
``Subsets with duplicates'' & Find subsets without creating duplicates & Backtracking (Sort + skip-equal-sibling) \\ \hline
\end{tabularx}

\section{Summary and Complexity Reference}

\noindent
\begin{tabularx}{\textwidth}{|l|l|l|X|}
\hline
\textbf{Task / Operation} & \textbf{Time} & \textbf{Space} & \textbf{Use When} \\ \hline
Factorial & $O(n)$ & $O(n)$ (call stack) & Computing $n!$ directly \\ \hline
Subsets & $O(n \cdot 2^n)$ & $O(n)$ (stack depth) & $n \le 20$ \\ \hline
Permutations & $O(n \cdot n!)$ & $O(n)$ & $n \le 8$ ($8! = 40320$) \\ \hline
Combinations ($n$ choose $k$) & $O(k \cdot C(n,k))$ & $O(k)$ & $n \le 20$ \\ \hline
Generate Parentheses & $O(n \cdot C_n)$ & $O(n)$ & $n \le 8$, $C_n$ is $n$-th Catalan number \\ \hline
N-Queens & $O(n!)$ worst case & $O(n)$ & $n \le 13$ practical \\ \hline
\end{tabularx}

\section{Glossary}

\subsection{Core Concepts}
\begin{itemize}
    \item \textbf{Recursion}: A function calling itself to solve a smaller piece of the same problem.
    \item \textbf{Base Case}: The condition that stops the recursion. Without it, the function calls itself forever.
    \item \textbf{Recursive Case}: The part of the function that breaks the problem down and makes the recursive call.
    \item \textbf{Call Stack}: The computer's internal memory structure (like a stack of sticky notes) that keeps track of active functions.
    \item \textbf{Stack Frame}: A single ``sticky note'' on the call stack containing a function's local variables.
    \item \textbf{Stack Overflow}: A fatal error that happens when recursion goes too deep and the computer runs out of stack frames.
    \item \textbf{Backtracking}: Exploring all possible options for a problem by trying a path and ``undoing'' the choice if it fails.
    \item \textbf{Decision Tree}: A visual way to map out all the choices made during a backtracking algorithm.
    \item \textbf{Search Space}: The total universe of all possible solutions you have to look through.
    \item \textbf{Pruning}: Stopping the exploration of a path early because you realize it cannot possibly lead to a valid answer.
    \item \textbf{Constraint}: A hard rule that a valid solution must follow.
    \item \textbf{Feasibility Check}: Code that tests whether making a choice violates a constraint.
    \item \textbf{Permutation}: An arrangement of items where the order matters (for example, [1,2] is different from [2,1]).
    \item \textbf{Combination}: A selection of items where the order does not matter.
    \item \textbf{Subset (Power Set)}: A set formed by taking any number of elements from the original set.
    \item \textbf{N-Queens}: A classic backtracking puzzle requiring you to place $N$ queens on an $N \times N$ chessboard safely.
\end{itemize}

\subsection{Mathematical Symbols}
\begin{itemize}
    \item \textbf{$n!$}: Factorial. The product of all positive integers up to $n$.
    \item \textbf{$C(n,k)$ or $\binom{n}{k}$}: ``N choose K''. The number of ways to pick $k$ items from a set of $n$ items.
    \item \textbf{$C_n$}: The $n$-th Catalan number. It represents the number of valid arrangements of brackets, among other things.
    \item \textbf{$2^n$}: 2 raised to the power of $n$. Represents the total number of subsets of an $n$-element set.
\end{itemize}

\section{Further Reading}

\subsection{Free Online Resources (Start Here)}
\begin{itemize}
    \item \textbf{cp-algorithms.com}: Excellent code-focused articles on combinatorial search and backtracking.
    \item \textbf{Art of Problem Solving (AoPS) Wiki}: Beginner-friendly math explanations on combinatorics and recursion.
    \item \textbf{Project Euler (projecteuler.net)}: Try solving problems that involve counting or generation (for example, Problem 15: lattice paths, Problem 24: lexicographic permutations).
    \item \textbf{Brilliant.org}: Offers interactive and visual courses on Recursion and Combinatorics.
\end{itemize}

\subsection{Books (When You Are Ready to Go Deeper)}
\begin{itemize}
    \item \textit{Competitive Programming 3} by Steven Halim. Covers recursive backtracking, pruning, and problem-solving strategies. (Intermediate difficulty, read when comfortable with basic implementations).
    \item \textit{Cracking the Coding Interview} by Gayle Laakmann McDowell. Chapter 8 covers recursion and dynamic programming. (Beginner-friendly, excellent for interview prep).
    \item \textit{The Algorithm Design Manual} by Steven Skiena. Chapter 7 covers combinatorial search and backtracking in depth. (Intermediate difficulty, strong practical insights).
    \item \textit{Introduction to Algorithms (CLRS)}. Contains rigorous chapters on recursive algorithms and the master theorem. (Advanced, graduate-level mathematical reference, read when formal analysis is required).
\end{itemize}

\section*{Version History}
\addcontentsline{toc}{section}{Version History}

\noindent
\begin{tabularx}{\textwidth}{|l|X|}
\hline
\textbf{Date} & \textbf{Change} \\ \hline
2026-10-03 & Document created. \\ \hline
2026-10-03 & Expanded backtracking traces, diagram trees, code standards, duplicate logic, and pattern rules. \\ \hline
\end{tabularx}

\end{document}
